\documentclass[conference,a4paper]{IEEEtran}
\usepackage[T5]{fontenc}
\usepackage[utf8]{inputenc}
\usepackage[vietnamese]{babel}
\usepackage{amsmath,amssymb,amsfonts}
\usepackage{algorithmic}
\usepackage{graphicx}
\usepackage{textcomp}
\usepackage{xcolor}
\usepackage{hyperref}
\usepackage{float}
\usepackage{booktabs}
\usepackage{multirow}
\usepackage{subcaption}
\usepackage{listings}
\usepackage{cite}
\usepackage{balance}

\lstset{
basicstyle=\ttfamily\footnotesize,
breaklines=true,
frame=single,
numbers=left,
numberstyle=\tiny\color{gray},
keywordstyle=\color{blue},
stringstyle=\color{red},
commentstyle=\color{green!50!black},
showstringspaces=false
}

\begin{document}

\title{Hệ thống Hỗ trợ Chẩn đoán Ung thư Vú Từ Hình Ảnh Y Khoa Ứng dụng Học Sâu (Deep Learning)\\
\vspace{0.2cm}
{\Large Báo cáo Đồ án Cuối kỳ Môn học Trí tuệ Nhân tạo}
}


\author{
\IEEEauthorblockN{Phạm Vũ Trần Anh, Đào Văn Huy Bình, Nguyễn Đức Dũng}
\IEEEauthorblockA{Nhóm 1, Lớp KHMT 18-01, Khoa Công Nghệ Thông Tin \\
Trường Đại Học Đại Nam, Việt Nam}
\and
\IEEEauthorblockN{ThS. Lê Trung Hiếu, KS. Nguyễn Thái Khánh}
\IEEEauthorblockA{Giảng viên hướng dẫn, Khoa Công Nghệ Thông Tin \\
Trường Đại Học Đại Nam, Việt Nam}
}

\maketitle


\begin{abstract}
Ung thư vú tiếp tục là nguyên nhân gây tử vong hàng đầu liên quan đến ung thư ở phụ nữ trên toàn thế giới. Chẩn đoán lâm sàng truyền thống phụ thuộc nhiều vào năng lực và kinh nghiệm của bác sĩ thông qua việc đọc ảnh nhũ ảnh, siêu âm và ảnh vi thể sinh thiết, dẫn đến tốn thời gian và mang tính chủ quan. Đồ án này trình bày toàn bộ quy trình thiết kế, phát triển và đánh giá một Hệ thống Hỗ trợ Chẩn đoán (Computer-Aided Diagnosis - CAD) sử dụng kỹ thuật Học Sâu (Deep Learning), đặc biệt là các biến thể Mạng nơ-ron Tích chập (CNN) hiện đại: EfficientNet-B3 và ResNet-50. Đối với phân loại ảnh siêu âm (Ultrasound), mô hình EfficientNet-B3 được áp dụng học chuyển giao (Transfer Learning) để nhận diện 3 phân lớp (Bình thường, Lành tính, Ác tính). Đối với dữ liệu mô bệnh học đa phân giải (Histopathology), mô hình ResNet-50 được sử dụng cho bài toán phân loại nhị phân mô ung thư ác tính. Bên cạnh việc tiền xử lý ảnh và tăng cường dữ liệu phức tạp để đối phó với hiện tượng mất cân bằng nhãn y khoa, hệ thống còn tích hợp công cụ trực quan hóa Grad-CAM để chiếu (heatmap) các khu vực có vết thương, cung cấp khả năng giải thích AI (XAI) nhằm tăng tính ứng dụng trong y tế lâm sàng. Toàn bộ các luồng suy diễn mô hình được phục vụ thông qua một backend Flask RESTful API và giao diện Web (Frontend) phản hồi thời gian thực. Đồ án không chỉ dừng lại ở chỉ số lý thuyết (F1-Score, ROC-AUC) mà còn hướng trực tiếp tới sản phẩm thực tiễn hóa sức khỏe số.
\end{abstract}

\begin{IEEEkeywords}
Ung thư vú, Học Sâu (Deep Learning), Convolutional Neural Networks (CNN), Hệ thống CAD, Phân tích Hình ảnh Y khoa, ResNet, EfficientNet, Flask, Grad-CAM.
\end{IEEEkeywords}

\section{Giới thiệu chung}
\subsection{Bối cảnh và Động lực}
Ung thư vú là một trong những loại ung thư với tỷ lệ mắc phải ngày càng trẻ hóa. Theo báo cáo thường niên của Tổ chức Y tế Thế giới (WHO), mỗi năm có hàng triệu ca mắc ung thư tuyến vú mới và sự sống còn của bệnh nhân phụ thuộc sống còn vào việc khối u được phát hiện ở giai đoạn nào. Nếu chẩn đoán ung thư ở giai đoạn sơ khởi (giai đoạn 0 hoặc 1), tỷ lệ sinh tồn trên 5 năm của bệnh nhân có thể đạt đến 99\%. Tuy nhiên, đối với bệnh nhân ở giai đoạn di căn hạch bạch huyết hoặc di căn xa, tỷ lệ này sụt giảm thê thảm. Do đó, khám sàng lọc định kỳ (Screening) được các tổ chức ung thư quốc tế đặc biệt khuyến nghị.

Hiện nay, các xét nghiệm hình ảnh thông dụng bao gồm Chụp X-quang vú (Mammography), Siêu âm vú (Breast Ultrasound) và Chụp cộng hưởng từ (Magnetic Resonance Imaging - MRI). Khi xuất hiện bất thường (ví dụ khối u không rõ rìa, có cấu trúc vi vôi hóa), sinh thiết tế bào khối u ra ngoài cơ thể và quan sát dưới kính hiển vi quang học (Histopathology) là tiêu chuẩn vàng để xác định bản chất (Lành tính - Benign hay Ác tính - Malignant).

Tuy nhiên, thủ tục phân tích y khoa truyền thống mắc phải các giới hạn sau:
\begin{enumerate}
\item \textbf{Phụ thuộc Chuyên gia:} Hình dáng khối u trên phim cực kì khó đọc, bị lẫn vào cấu trúc mô tuyến (density). Ngay cả các bác sĩ X-quang tay nghề cao cũng phải đối mặt với sai sót chẩn đoán (âm tính giả).
\item \textbf{Nhiễu loạn Thiết bị:} Máy siêu âm sinh ra nhiễu hạt (speckle noise) do sóng âm phân tán, có thể vô tình làm lu mờ đường biên mô u ác.
\item \textbf{Khối lượng công việc khổng lồ:} Bác sĩ phòng khám tuyến đầu luôn trong tình trạng quá tải (Burnout), khiến khả năng bỏ sót tổn thương trong lúc mệt mỏi gia tăng.
\end{enumerate}

Với sự phát triển bùng nổ của Trí tuệ Nhân tạo (AI), hệ thống CAD (Computer-Aided Diagnosis) đóng vai trò như một màng lọc tự động. Một trợ lý AI không biết mệt mỏi sẽ tiếp nhận ảnh trước, phân tích với tốc độ mili-giây, khoanh vùng khu vực nguy hiểm và đặt sẵn cảnh báo xác suất ác tính. Bác sĩ chuyên khoa sẽ nhìn vào dự đoán này để có sự đối chiếu (Second Opinion), thay vì phải quét trơn mắt như trước đây, giảm mệt mỏi và nâng cao hiệu quả y tế.

\subsection{Bài toán và Phát biểu Yêu cầu}
Mục tiêu bao quát của đồ án môn học này là xây dựng một "Hệ thống chuyên gia nội bộ" dùng thuật toán học sâu - có khả năng nhận các dòng bit dữ liệu ảnh bệnh nhân truyền qua HTTP/REST, xử lý nó xuyên suốt cấu trúc mạng nơ-ron hằng chục triệu tham số, và xuất ra báo cáo số học về xác suất bệnh tính.

Đồ án được phân rã thành các bài toán thực thi như sau:
\begin{itemize}
\item \textbf{Xử lý nhiễu và Tối ưu dữ liệu đầu vào:} Xử lý tập ảnh thu thập từ Internet (BUSI, BreaKHis), thay đổi không gian màu, cắt ảnh (Center Crop), làm sắc nét, và thực hiện Data Augmentation ngẫu nhiên nhắm tới làm giàu tập huấn luyện (Tập Train) nhằm ngăn biến dạng ngẫu nhiên.
\item \textbf{Mô hình hóa (Modeling):} Nghiên cứu, thiết kế hàm mất mát (Loss Function), khởi tạo siêu tham số (Hyperparameters), và chuyển giao tri thức từ cơ sở dữ liệu vạn vật ImageNet (1.2 triệu ảnh) phục vụ tinh chỉnh (fine-tuning) các lớp tuyến tính cuối vào ma trận không gian y học tuyến vú.
\item \textbf{Đóng gói thành Sản phẩm Công nghệ (Productization):} Dịch mã thuần Deep Learning sang một Server phục vụ công cộng. Đảm bảo luồng xử lý đồng thời qua API chuẩn xác. Web tĩnh được triển khai cho bác sĩ sử dụng.
\end{itemize}

\subsection{Ý nghĩa Khoa học và Thực tiễn}
Đóng góp chính của báo cáo tập trung vào:
\begin{itemize}
\item Khảo sát chi tiết hai luồng dữ liệu y khoa dị biệt (siêu âm phản xạ sóng âm và sinh thiết soi màu quang phân giải cực cao).
\item Tối ưu hóa mô hình học sâu CNN để cân nhắc tỉ lệ trade-off giữa năng lực nhận thức (Accuracy/Recall) và Tốc độ dự đoán (Inference Latency).
\item Mã nguồn mã nguồn mở (Open-source implementation) dễ cài lại và ứng dụng.
\end{itemize}

Bố cục phần còn lại: Phần II trình bày Tổng quan các nghiên cứu liên quan. Phần III đi sâu vào cơ sở lý thuyết kiến trúc mạng Nơ-ron. Phần IV giải bày toàn bộ quy trình thiết kế dữ liệu. Phần V là huấn luyện mô hình học. Phần VI là đánh giá và Các độ đo thực nghiệm. Phần VII đưa vào giải thích mô hình AI. Cuối cùng, Phần VIII và IX bàn về hệ thống công nghệ và tầm nhìn.

\section{Tổng quan các nghiên cứu liên quan (Literature Review)}
\subsection{Các phương pháp Machine Learning Truyền thống}
Trước giai đoạn những năm 2012 (trước khi AlexNet bùng nổ), CAD Y khoa sống dựa vào Computer Vision truyền thống. Các nhà nghiên cứu phải tự tay (hand-craft) tạo ra bộ trích xuất toán học.
Ví dụ, để bác sĩ máy hiểu được thế nào là khối u, các phương pháp cũ thường trích xuất đặc trưng hình thái học (Morphological Features - chu vi, diện tích u, độ xù xì rìa) và đặc trưng dệt (Textural Features - ma trận đồng hiện mức xám GLCM, GLRLM). Sau khi các biến số này được chuyển thành ma trận đặc trưng $1D$, thuật toán học máy cổ điển như Máy Vector Hỗ Trợ (Support Vector Machine - SVM), Random Forest (RF) hay K-Láng giềng gần nhất (KNN) sẽ tìm ra các đường siêu phẳng phân loại. Hạn chế chết người của phương pháp này là nó bị phụ thuộc lớn vào việc khu vực khối u phải được "Cắt bóng" (Segment) thủ công thật chính xác từ bác sĩ rồi nhập vào máy, do đó chưa có khái niệm Fully-Automated (Hoàn toàn tự động).

\subsection{Các phương pháp Deep Learning}
Với sự có mặt của GPGPU (Graphics Processing Unit tính toán phổ quát), Mạng nơ-ron Tích chập (CNN) đã học được cách tự trích xuất đặc trưng hình ảnh với các lớp chập (Convolution Layers) mà không cần bóc tách y khoa thủ công.
\begin{itemize}
\item \textbf{Spanhol et al. (2016)} \cite{breakhis} thu thập tập BreaKHis và thử nghiệm bằng AlexNet. Tuy vậy vì AlexNet khá mỏng (nông) và không có công cụ điều chỉnh batch, kết quả vẫn gặp tình trạng Overfitting mạnh.
\item Nối tiếp theo đó, \textbf{VGG-16} hay \textbf{InceptionNet} ra đời mang lại Recall cao hơn nhưng vì cấu trúc số lượng weights khổng lồ, khó ứng dụng trên các máy tính y tế có phần cứng thấp (như trạm siêu âm cũ).
\item \textbf{Transfer Learning} trở thành tiêu chuẩn vàng. Nhà nghiên cứu Al-Dhabyani \cite{busi} xuất bản tập ảnh siêu âm BUSI vào năm 2020. Các kiến trúc Residual như \textbf{ResNet} được ứng dụng rộng rãi. ResNet-50 học cực kỳ nhạy bén các bất thường tế bào, tuy vậy với ảnh siêu âm nhiễu, Residual blocks thường học nhầm noise âm thành khối u.
\item Để vượt rào, kiến trúc \textbf{EfficientNet} (Tan và Le, Google Brain) kết hợp tỷ lệ cân bằng độ nông-sâu của không gian mạng. Do đó Đồ án của chúng tôi đã quyết định kết hợp EfficientNet cho ảnh ngoại khoa (siêu âm nhiễu) và ResNet cho ảnh tế bào chi tiết nội khoa sinh thiết để đạt được độ chính xác bù trừ tuyệt đối nhất.
\end{itemize}

\section{Cơ sở Lý thuyết}
Máy tính chỉ hiểu hình ảnh dưới dạng cấu trúc đa chiều $H \times W \times C$ (Chiều Cao $\times$ Chiều rộng $\times$ Số lượng Kênh RGB cục bộ). Khác với mạng perceptron nhiều lớp truyền thống (MLP) sẽ trải phẳng toàn bộ ảnh đầu vào ra $1D$, dẫn đến làm vỡ cấu trúc tương quan địa lý giữa các điểm ảnh liền kề, CNN xử lý theo mô hình trượt lọc.

\subsection{Lớp Tích chập (Convolutional Layer) và Toán học}
Trong lớp này, ma trận đặc trưng biến dổi bằng cách quét các "kernel" (hoặc filter) khối lượng $N \times N$ khắp tấm ảnh với một bước nhảy (Stride - $s$) cố định.
Xét ma trận ảnh 2 chiều đầu vào là $I$, biến đổi hàm bộ lọc không gian $K$:
\begin{equation}
O(i, j) = (I * K)(i, j) = \sum_{m=0}^{M-1} \sum_{n=0}^{N-1} I(i \cdot s + m, j \cdot s + n) K(m, n)
\end{equation}
Tham số $K(m,n)$ chính là "tri thức" (weights) mà hệ thống tự động cập nhật trong lúc train. Mạng chập học được các cạnh (Edges), cấu trúc vòng, đốm sẫm ở các tầng chập đầu tiên và tổ hợp thành hình thái khối u lớn ở các tầng trong sâu nhất.
Để ngăn việc ma trận thu gọn liên tục qua từng đợt cuộn, phương pháp viền biên đệm giá trị số 0 (Zero-Padding) được ứng dụng đồng đều, biểu thị là $p$.
\begin{equation}
H_{out} = \lfloor \frac{H_{in} + 2p - M}{s} \rfloor + 1
\end{equation}

\subsection{Hàm Kích hoạt Phi tuyến tính (Activation Function)}
Nếu chỉ tích chập nối tiến, mạng bản chất sẽ sụp về một phương trình tuyến tính đại số duy nhất $y = Ax+B$, không thể giải nổi những biểu đồ không gian khuyết tật khối u vốn phi tuyến.
\begin{itemize}
\item \textbf{ReLU (Rectified Linear Unit)}: $f(x) = \max(0, x)$. Đơn giản và loại bỏ vanishing gradient phía số dương.
\item \textbf{Swish}: Được tìm ra tự động do thuật toán AutoML của tập đoàn Google. Swish vượt trội hơn ReLU ở mô hình càng sâu, đây cũng là hàm nền tảng kích hoạt của hệ EfficientNet:
\begin{equation}
f(x) = x \cdot \text{sigmoid}(\beta x) = \frac{x}{1 + e^{-\beta x}}
\end{equation}
Đạo hàm của Swish luôn mượt mà chứ không chứa góc gãy ngay điểm 0 như ReLU.
\end{itemize}

\subsection{Lớp gộp (Pooling) và Lớp kết nối đầy đủ (FC)}
- \textbf{Max Pooling:} Giảm kích thước không gian bằng cách lấy trị giá điểm ảnh cao độ nhất (đặc trưng chói quan trọng nhất) trong khu vực trượt $(2 \times 2)$. Điều này đem lại tính "Bất biến cục bộ tĩnh đối với phép tịnh tiến". Dù khối u có lệch đi đôi chút trên ảnh, nó vẫn rơi vào khoảng Max Pooling và máy luôn nhận ra.
- \textbf{Global Average Pooling (GAP):} Ở tầng cuối thay vì lặp flatten thành $1D$ Vector (tốn cả triệu tham số kết nối với lớp Dense theo ma trận chéo), GAP tính trung bình mọi thông số trên 1 channel map còn lại duy nhất một con số. Điều này tránh được Overfitting cho y khoa cực kỳ hữu hiệu.

Đầu ra cùng cùng đẩy vào Lớp \textbf{Softmax} đa phân loại:
\begin{equation}
\text{Softmax}(\mathbf{z})_i = \frac{e^{z_i}}{\sum_{j=1}^{K} e^{z_j}}
\end{equation}
với $\mathbf{z}$ là chuỗi kết quả tuyến tính (logits) thô. Mảng xuất này biến tổng xác suất mọi nhãn thành $1.0$ ($100\%$).


\subsection{ResNet-50 (Residual Network)}
Do giới hạn của cấu trúc nơ-ron, khi xếp chồng đến tầng thứ 20, các tham số mất tín hiệu đạo hàm lùi dẫn đến sai số lan truyền, hiện tượng suy thoái chất lượng gradient biến mất. Dr. Kaiming He \cite{resnet} xuất bản kỹ thuật Skip Connection.
Nếu $\mathbf{x}$ là đặc trưng đầu vào cùa khối nơ ron, ma trận cần học sẽ là hàm Dư:
\begin{equation}
F(\mathbf{x}) = \mathcal{H}(\mathbf{x}) - \mathbf{x}
\end{equation}
Đầu ra sẽ ép khối đi qua hàm cộng ngang (Identity Shortcut): $y = \text{ReLU}(F(\mathbf{x}) + \mathbf{x})$.
Mô hình Resnet-50 ở đồ án sử dụng cơ chế Bottleneck (Trục xoắn): giảm số chiề channel từ 256 xuống 64 với lớp chập $(1 \times 1)$, xử lý với $(3 \times 3)$, và giải nén channel lại 256 dùng $(1 \times 1)$. Cách làm tiết kiệm khối lượng FLOPS một cách đáng kể. Cơ chế này giải quyết triệt để sự cố suy giảm đạo hàm (Vanishing Gradient) khi các lớp sâu hơn không học được bất kỳ đặc trưng nào do độ dốc tiệm cận về 0 sau chuỗi phép nhân liên tiếp. Áp dụng vào miền hình ảnh sinh thiết y khoa (BreaKHis), các Residual Block cấp cao đóng vai trò sinh tử trong việc bảo lưu nguyên vẹn hình thái lớp màng và nhân tế bào từ các tầng nông truyền nguyên khối xuống tầng sâu mà không bị làm mờ bởi hàng chục lần nhân ma trận.

\subsection{EfficientNet và Compound Scaling}
Mạng CNN truyền thống thường chỉ "cơi nới" mạng theo một chiều sâu (ResNet), chiều rộng (WideResNet) hoặc đẩy độ phân giải. EfficientNet \cite{efficientnet} toán học hóa quy trình này. Ký hiệu mạng là hàm $\mathcal{N} = \bigodot_{i=1...s} \mathcal{F}_i^{L_i} (X_{\langle H_i, W_i, C_i \rangle})$. Bài toán của máy là tối ưu Accuracy với ràng buộc tài nguyên $Memory$ cố định. Phương trình hệ số phức hợp $\phi$:
\begin{eqnarray}
\text{depth (d)} = \alpha^\phi \\
\text{width (w)} = \beta^\phi \\
\text{resolution (r)} = \gamma^\phi \\
s.t. \quad \alpha \cdot \beta^2 \cdot \gamma^2 \approx 2
\end{eqnarray}
EfficientNet-B3 trong backend API được sử dụng khối lượng tham số xấp xỉ 12 triệu parameter (Cực nhẹ nhẹ) nhưng đánh bại các mạng khổng lồ 100 triệu tham số trong tập ImageNet. Khối trung tâm làm nên sức mạnh mạng này là \textbf{MBConv (Mobile Inverted Bottleneck Convolution)}, kèm theo cơ chế ép chắt không gian \textbf{Squeeze-and-Excitation (SE)}. Đối với hiện trạng nhiễu hạt bụi \textit{speckle} dày đặc trên phim siêu âm tuyến vú, cơ chế chú ý SE đóng vai trò như một màng lọc tự nhiên: Nó tự học cách "vặn nhỏ Volume" của các kênh chứa rác nhiễu (Background Noise) và bộc lộ các kênh trích xuất mầm gai góc của khối u ác tính. Sự kết hợp giữa khả năng nới rộng đa chiều (Compound Scaling) và MBConv biến EfficientNet thành kiến trúc lý tưởng nhất phục vụ đồ án (Tỷ lệ Trade-off giữa độ chuẩn xác và tốc độ API thời gian thực cực kỳ tuyệt vời).

\section{Dữ liệu và Phương pháp Đề xuất Đầu vào}
Một pipeline máy học chỉ thông minh ngang với lượng dữ liệu mà kỹ sư bơm vào (Garbage in, garbage out).

\subsection{Tập dữ liệu Siêu âm Vú (BUSI)}
Được thu thập vào năm 2018 tại Bệnh viện Baheya, Cairo, Ai Cập \cite{busi}. Chúng bao gồm thông số thực y khoa của 600 bệnh nhân. Các chẩn đoán được xác nhận bằng sinh thiết tế bào học thật và phân chia:
\begin{table}[H]
\centering
\caption{Thống kê định lượng nhãn BUSI Ultrasound Dataset}
\begin{tabular}{lcc}
\toprule
\textbf{Nhãn lớp Y khoa} & \textbf{Số Lượng Ảnh gốc} & \textbf{Kết quả mong muốn} \\
\midrule
Bình Thường (Normal) & 133 & Âm tính hoàn toàn \\
Lành Tính (Benign) & 437 & Âm tính ác (Khối sợi) \\
Ác Tính (Malignant) & 210 & Dương tính ác tính \\
\midrule
\textbf{Tổng cộng} & 780 & 3 Lớp nhãn phân biệt \\
\bottomrule
\end{tabular}
\end{table}
Tỷ lệ chênh lệch là khổng lồ (Lành tính chiếm gần gập 3.5 lần nhãn Bình thường). Không được phép để AI học lố theo nhãn đa số.

\begin{figure}[H]
\centering
\begin{minipage}{0.45\linewidth}
\centering
\includegraphics[width=\linewidth]{benign (4).png}
\centerline{(a) Mẫu Lành tính (Benign)}
\end{minipage}
\hfill
\begin{minipage}{0.45\linewidth}
\centering
\includegraphics[width=\linewidth]{malignant (2).png}
\centerline{(b) Mẫu Ác tính (Malignant)}
\end{minipage}
\vspace{0.2cm}
\caption{Hình ảnh minh họa trích xuất thực tế từ bộ dữ liệu Siêu âm Vú (BUSI). Khối u ác tính thường có viền gai góc, ranh giới không rõ ràng và tối màu (Hypoechoic) hơn so với nốt sợi tuyến vú lành tính.}
\end{figure}

\subsection{Tập dữ liệu Vi Thể Sinh Thiết (BreaKHis)}
Các mẫu khối u được nhuộm bởi thuốc nhuộm mô cơ bản Hematoxylin \& Eosin (H\&E). Hematoxylin nhuộm nhân tế bào màu xanh đen-tím, trong khi Eosin nhuộm chất nguyên sinh và các sợi liên kết màu hồng và đỏ. Mạng học sâu sẽ dựa theo sự biến dị của tỉ lệ viền nhân \cite{breakhis} để dự đoán.
Tổng 7.909 hình ảnh phóng đại phân giải cao. Tuy nhiên, thay vì phân tích đơn hình (Single patch), đồ án đã downscale các ảnh này về chung chiều $224 \times 224$ vector.
\begin{itemize}
\item Lành tính sinh thiết (Benign - Adenosis, Fibroadenoma): Tổng lượng 2,480 files đại diện. Về mặt sinh lý, đây là những quần thể tế bào loạn sản cấu trúc sợi tuyến vú nhưng hoàn toàn không có dấu hiệu thâm nhiễm dính màng. Quỹ đạo nhân tế bào của chúng vẫn giữ được tỷ lệ diện tích bảo toàn tương xứng với tương bào bên ngoài (Nucleus-to-Cytoplasm Ratio ổn định).
\item Ác tính sinh thiết (Malignant - Carcinoma mảng biểu mô, di căn nội ống NST): Khối lượng 5,429 files đặc cách. Tổ chức hiển vi của những mô này đặc thủng bởi mật độ tế bào dày đặc che kín nền đen của mô, tế bào lộn xộn mất phương hướng. Đặc biệt, nhân tế bào có hiện tượng phân bào nhân chia cực kỳ dị dạng (Atypical Mitosis), biến dạng phồng lên và đổi màu tối đậm đặc (Hyperchromasia), phá vỡ toàn bộ cấu trúc kiến trúc màng đáy niêm mạc xung quanh để tràn vào mao mạch tản ra toàn cơ thể.
\end{itemize}

Sự khác biệt mang tính chất sinh hóa di truyền học này cực kỳ khắc nghiệt và khó nắm bắt với mắt thường của các bác sĩ giải phẫu bệnh học non trẻ (Junior Pathologists). Đặc biệt là dưới sức ép hàng núi bệnh án mỗi ngày, sự mệt mỏi thoái hóa điểm vàng võng mạc qua nhiều giờ dán mắt trực tiếp vào ống kính hiển vi quang học luôn trực chờ kéo theo tỷ lệ sinh lỗi cao (False Negative rate). Chuyên gia y khoa có thể mệt, nhưng mã nguồn nhị phân Toán học thì không. Đó chính là hệ tầng sinh thái hoàn hảo tuyệt đối nhằm giải phóng thuật toán Tích chập Đa chiều Convolutional Neural Network (mảnh ghép sinh ra đặc quyền riêng để nhận diện Textures và Cạnh Viền Hình Học) vươn lên giành lấy vương miện lợi thế phân loại cao nhất.

\subsection{Tiền Xử lý (Preprocessing) Cốt lõi}
Hình dạng bức ảnh từ file dạng thô (.png, .jpg, .dicom đa cực) bị chuyển vị về lưới PyTorch Tensor là $(C \times H \times W)$ (Channel, Height, Width).
\begin{equation}
x_{normalized} = \frac{x - \mu}{\sigma}
\end{equation}
Trong đó vector ma trận hằng số chuẩn hóa $\mu = [0.485, 0.456, 0.406]$ là trung bình, và $\sigma = [0.229, 0.224, 0.225]$ là độ lệch chuẩn màu đã được đúc kết từ hàng tỷ lượt xử lý từ máy chủ Google.
\textbf{Tăng cường dữ liệu thời gian thực (On-the-fly Data Augmentation)}:
Nhằm phòng tránh mô hình "Học Thuộc Lòng" - Overfitting (Khi Training Loss giảm về 0 nhưng Testing Accuracy lại loanh quanh 50\%), mỗi batch mini 16 hình ảnh đều được "gây nhiễu có chủ đích". Code snippet ở khối `utils.py` backend:
\begin{lstlisting}[language=Python]
import torchvision.transforms as transforms
train_transforms = transforms.Compose([
transforms.Resize((256, 256)),
transforms.RandomRotation((-15, 15)), # Xoay hinh 15 do
transforms.RandomHorizontalFlip(p=0.5), # Lat hinh
transforms.CenterCrop(224), # Cat khung luoi tieu chuan
transforms.ToTensor(),
transforms.Normalize(...)
])
\end{lstlisting}

Phép xoay tịnh tiến ma trận biến đổi affine được định nghĩa:
\begin{equation}
\begin{bmatrix} x' \\ y' \end{bmatrix} =
\begin{bmatrix} \cos\theta & -\sin\theta \\ \sin\theta & \cos\theta \end{bmatrix}
\begin{bmatrix} x - x_c \\ y - y_c \end{bmatrix} + \begin{bmatrix} x_c \\ y_c \end{bmatrix}
\end{equation}

Việc lật Horizontal bảo đảm máy không học "Khối u ác tính lúc nào cũng xuất hiện nằm bên tay trái". Khối u hoàn toàn đối xứng trong mô vú, bác sĩ cầm đầu dò siêu âm (Trandsducer) đổi góc độ là hình ảnh bị lật theo. Do đó Augmentation là một kỹ thuật ép buộc sự bất biến (Invariance) hình học.

Bên cạnh đó, vì kích thước ảnh đàu vào thô do siêu âm và máy quang học chụp ra không hề đồng nhất, khâu \textbf{Resize Interpolation (Nội suy Song tuyến tính)} được áp dụng bắt buộc để định hình tensor về ma trận vuông tỷ lệ $256 \times 256$ trước khi thuật toán RandomCrop cắt chính xác vùng trung tâm $224 \times 224$ đưa vào mạng nơ-ron. Phương pháp Bilinear tính toán sắc độ điểm ảnh mới dựa trên trung bình gia quyền của 4 điểm ảnh lân cận nhất trong ảnh gốc đại số:
\begin{equation}
\resizebox{0.95\linewidth}{!}{$
f(x,y) \approx \frac{1}{(x_2-x_1)(y_2-y_1)} [x_2-x \quad x-x_1] \begin{bmatrix} f(Q_{11}) & f(Q_{12}) \\ f(Q_{21}) & f(Q_{22}) \end{bmatrix} \begin{bmatrix} y_2-y \\ y-y_1 \end{bmatrix}
$}
\end{equation}
Sự kết hợp đa dạng các phép biến đổi này giúp đồ án khắc phục triệt để bài toán thiếu hụt mẫu ung thư nghiêm trọng (Chỉ có vỏn vẹn 210 ảnh ác tính nguyên bản ở tập BUSI), từ đó khuếch đại không gian vector mẫu lên gấp hàng chục lần mà không làm mất đi tính biểu trưng nguyên bản của y khoa tế bào.
\section{Huấn luyện AI và Tối ưu Hóa (Training Flow)}
Quy trình được điều phối tự động hoàn toàn bằng tập lệnh Python. Tập Validation được Random Split giữ ngẫu nhiên mầm cố định $20\%$ so với tổng số lượng, không bao giờ được phép tham gia Backpropagation cập nhật Gradient, giúp kỹ sư giám sát độc lập tính khách quan học quá đà của trí tuệ nhân tạo.

\subsection{Cấu hình Môi trường Phần cứng và Siêu tham số (Hyperparameters)}
Việc huấn luyện các kiến trúc CNN sâu nặng nề tốn GPU (Heavy-weight networks) yêu cầu tài nguyên phần cứng tính toán song song (Parallel Computing) cao độ. Hệ thống phát triển được chuẩn hóa đồng nhất trên thư viện lõi PyTorch, kết hợp nhân tính toán lớp chập ma trận GPU NVIDIA (CUDA Device) để tối ưu hóa triệt để thuật toán Tensor trôi với độ trễ tối thiểu (Low Latency).

\begin{table*}[t]
\centering
\caption{Bảng Đặc tả Môi trường Huấn luyện máy học và cấu hình Siêu Tham số}
\renewcommand{\arraystretch}{1.5}
\resizebox{\textwidth}{!}{
\begin{tabular}{|l|l|}
\hline
\textbf{Thành phần / Tham số} & \textbf{Thiết lập Cấu hình Thực nghiệm Đầu máy} \\ \hline
Nền tảng Tăng tốc (Hardware) & Nvidia Geforce RTX 3000 Series (hoặc Tiêu chuẩn ngang) \\ \hline
Kiến trúc Thư viện Trí Tuệ Deep Learning & PyTorch Framework (Backend Toán học CUDA) \\ \hline
Giải thuật Tối ưu Trọng số (Optimizer) & \textbf{Adam} (Thuật toán Adaptive Moment Estimation) \\ \hline
Kích thước Lô Tải Lên RAM (Batch Size) & $16$ Hình ảnh khối/1 Lô cập nhật chênh lệch Tensor \\ \hline
Tốc độ Buớc Nhảy (Learning Rate $\alpha$) & $1 \times 10^{-4}$ (Kèm bộ giảm \textit{StepLR} $0.1$ mỗi chu kỳ 7) \\ \hline
Tổng Số Chu kỳ Lặp dữ liệu (Epochs MAX) & Lớn nhất $50$ Epochs (Kèm bộ ngắt sớm Early Stopping) \\ \hline
Hàm Lỗi Đánh giá Phương Sai (Loss) & Cross-Entropy Error đa phân loại chuẩn hóa \\ \hline
Thuật toán Luồng phân bổ GPU I/O & \texttt{torch.utils.data.DataLoader(num\_workers=4)} \\ \hline
\end{tabular}
}
\end{table*}

\subsection{Kiến trúc cấp phát dữ liệu (Dataset \& DataLoader Architecture)}
Trong Pytorch, việc nạp cùng lúc vài chục ngàn bức ảnh ($224 \times 224$ ma trận màu 3 Lớp) vào trực tiếp lõi RAM (Lên đến hàng chục Gigabytes) là một điều siêu việt không tưởng và không tối ưu. Kiến trúc hệ thống lập trình tuân thủ nghiêm ngặt chuẩn thiết kế nạp bộ nhớ vi mô (Micro-Batching Loading). Kế hoạch thiết kế mã nguồn chia thành:
\begin{itemize}
\item \texttt{ImageFolder Dataset}: Hàm tự động duyệt qua cây thư mục vật lý ổ cứng Của Bác Sĩ (File tree) để tự động đặt nhãn. Ví dụ: Thư mục có đuôi `malignant` thì mọi ảnh Load lên sẽ tự ép vào Target Vector label 01. Nhờ đó, người dùng thêm ảnh bệnh mới chỉ cần copy-paste paste không cần sửa một dòng code nào.
\item \texttt{DataLoader Worker}: Lớp xử lý điều phối trung chuyển. Khi Card màn hình GPU bắt đầu huấn luyện mẻ (batch) ảnh thứ $(K)$, CPU không chịu ngồi chơi đợi. Tùy chọn \texttt{num\_workers=4} ra lệnh cho Chip CPU khai thác 4 tiểu trình Thread con, đi trước GPU ngầm mở SSD chuẩn bị giải mã tập (batch) thứ $(K+1)$ và $(K+2)$ lên RAM sẵn. Kỹ thuật Lập trình đa luồng (Multi-processing Pipelining) này tránh hiện tượng Thắt cổ chai vòng cổ (I/O Bottleneck Barrier).
\end{itemize}
\textbf{Chiến lược Học Chuyển Giao (Transfer Learning) Thực Tiễn}: Codebase hệ thống không khởi tạo mô hình ngẫu nhiên (Training From Scratch) mà thực hiện lệnh tải trước Model Trọng số (Pre-trained Weights) đã tinh luyện tới kỳ cuối cùng từ tập \textit{ImageNet} (1.2 triệu tấm ảnh quy mô lớn). Khối lượng bộ nhớ nhận thức hình họa thô (Góc cạnh, vệt, màng vòng) được "Đóng băng" (Freeze Architecture) ở các Tầng Phân Lớp Lõi dưới. Kỹ sư đồ án chỉ phá hủy phân khu Classifier Base Layer cuối cùng và thay thế chúng thành $M=3$ Node Tuyến tính Đầu ra (Mô hình Siêu âm) hoặc $M=2$ (Mô hình Sinh thiết nhị phân). Việc kết hợp tái rèn giũa bằng giải thuật Tối ưu Hóa thay thế toàn bộ quá trình mò mẫm tham số cơ bản, đẩy nhanh tốc độ hội tụ siêu dẫn trên hạ tầng phần mềm chỉ trong vòng 25-50 Chu kỳ huấn luyện (Epochs).

\subsection{Hàm Mất mát (Loss Fuction)}
Hàm Multi-class Cross-Entropy Error. Gọi khối lượng dữ liệu train từng chu kỳ mini-batch là $\mathcal{B}$, với $M$ là số lượng nhãn kết quả (M=3 với siêu âm, M=2 với sinh thiết):
\begin{equation}
\mathcal{J}_{CE}(\mathbf{W}) = -\frac{1}{|\mathcal{B}|} \sum_{i=1}^{|\mathcal{B}|} \sum_{c=1}^{M} \mathbf{y}_{i, c} \log(\hat{\mathbf{y}}_{i, c}(\mathbf{W}))
\end{equation}
Trong đó $\mathbf{W}$ đại biểu cho tensor Trọng Số (Weights) và các Siêu Hệ số nghiêng (Biases) trong các Neuron.
Với tập độ lệch (Imbalance Data), những phiên bản thử nghiệm có thay thế CrossEntropy bằng Focal Loss (Giảm trọng số các nhãn học dễ - Down-weighting easy examples), nhưng trong Đồ án này chỉ tập trung bằng Data Augmentation đã đủ độ bao phủ bề mặt lõi.

\subsection{Giải thuật Gradient Optimizer Momentum - ADAM}
Một Optimizer thuần thô sơ như Gradient Descent thường (SGD - Sotchastic Gradient Descent) sẽ tự bị mắc kẹt tại những vực dốc địa phương (Local Minima Điểm yên ngựa) trong trường dữ liệu cong.
Hệ thống sử dụng Kỹ Thuật Adam (Adaptive Moment Estimation) thay thế. Tại mỗi bước lặp tham số $t$:
Đạo hàm trọng số $g_t = \nabla_\theta \mathcal{J}(\theta_{t-1})$.
Tích lũy moment bậc một (Vận tốc tính hướng) và moment bậc 2 (RMS Tính quy mô hướng):
\begin{eqnarray}
m_t = \beta_1 m_{t-1} + (1 - \beta_1) g_t \\
v_t = \beta_2 v_{t-1} + (1 - \beta_2) g_t^2
\end{eqnarray}
Áp dụng sửa thiên lệch (Bias Correction) do khởi tạo tại số 0: $\hat{m}_t = m_t / (1 - \beta_1^t)$ và $\hat{v}_t = v_t / (1 - \beta_2^t)$.
Cuối cùng bộ nơ-ron tự cập nhập tham số:
\begin{equation}
\theta_t = \theta_{t-1} - \frac{\alpha}{\sqrt{\hat{v}_t} + \epsilon} \hat{m}_t
\end{equation}
Với $\alpha \ll 0.0001$ là Learning Rate sơ khởi. $\beta_1=0.9, \beta_2=0.999, \epsilon=10^{-8}$. Các công thức này được Framework PyTorch tự động đóng gói bên trong module `torch.optim.Adam`.

\subsection{Kiểm soát suy giảm tốc độ học và Early Stopping}
Để tránh Learning Rate quá cao làm mạng vượt lố điểm vàng Minimum, Scheduler Learning Rate được cài đặt. \texttt{StepLR(step\_size=7, gamma=0.1)} sẽ chủ động nhân $0.1$ (giảm $10$ lần) kích thước bước học mỗi khi đi qua 7 Epochs liên tiếp.
Vòng lập giám sát Validation, nếu Error không giảm thêm dù đã qua nhiều chu kỳ (Patience = 10), hệ thống kích hoạt khóa phanh (Early Stopping) và trích đoạn bộ khung trọng số (Model Weights Save) đang đạt chất lượng cực điểm xuống ổ cứng (Định dạng Tensor file `.pth`).

\section{Kết quả Thực nghiệm Phân Tích (Evaluation Results)}
Mọi thuật toán tính được đều phải dựa vào Bảng chân trị thực tế Ma trận nhầm lẫn y khoa (Confusion Matrix Analysis).
\begin{itemize}
\item \textbf{True Positive (TP)}: Bệnh nhân bị Ung thư mảng $\rightarrow$ AI đọc Ung thư mảng. (Thành Công sinh mệnh).
\item \textbf{True Negative (TN)}: Người khỏe mạnh $\rightarrow$ AI đọc là bình thường (Giảm chi phí bệnh án không cần thiết).
\item \textbf{False Positive (FP)}: Người khỏe mạnh $\rightarrow$ AI đo là có Cục U (Gây hoang mang, cần phải siêu âm/sinh thiết lại một lần nữa tốn kém thời gian).
\item \textbf{False Negative (FN)}: \textbf{Bệnh nhân bị Bệnh $\rightarrow$ AI đọc Không sao thả về nhà}. Đây là hiện tượng thảm khốc lâm sàng (Fatal Clinical Error) khiến tế bào di căn và phá hoại tính mạng bệnh nhân vì nhầm tin.
\end{itemize}

Sứ mệnh tối thượng đồ án AI Y Sinh học là bẻ gãy mọi chỉ số chui lọt FN về Mức Zê-rô.
Độ nhạy (Sensitivity - Recall) được giám sát gắt gao.

\begin{equation}
\text{Sensitivy/Recall} = \frac{TP}{TP + FN}
\end{equation}

Ngoài ra, code hệ thống đánh giá tự động sinh báo cáo Confusion Matrix với hàng loạt các độ đo (Metrics) khắt khe khác:
\begin{itemize}
\item \textbf{Precision (Độ chuẩn xác)}: Tỉ lệ chẩn đoán phát hiện Bệnh thật sự trên tổng số các ca AI gióng lên chuông cảnh báo là CÓ Bệnh định hướng cho y bác sĩ:
\begin{equation} \text{Precision} = \frac{TP}{TP + FP} \end{equation}
\item \textbf{F1-Score (Trung bình Điều hòa)}: Là tổ hợp Harmonic Mean giữa Precision và Recall. Đối với tập dữ liệu bất cân xứng khổng lồ, F1-Score là thước đo đánh giá thực lực vàng của mạng nơ-ron:
\begin{equation} F1 = 2 \times \frac{\text{Precision} \times \text{Recall}}{\text{Precision} + \text{Recall}} \end{equation}
\item \textbf{Accuracy (Độ chính xác Ròng đo trên Khung chuẩn)}: Tỉ lệ dự đoán khớp đúng mọi mặt trận lâm sàng bệnh thực thụ và người bình thường tổng:
\begin{equation} \text{Accuracy} = \frac{TP + TN}{TP + TN + FP + FN} \end{equation}
\end{itemize}

\subsection{Kết quả Thử nghiệm với EfficientNet-B3 (Siêu Âm)}
Với bộ Test độc lập (Khoảng 156 bệnh nhân):
\begin{table}[H]
\centering
\caption{Chỉ số Đánh giá của Mạng Siêu Âm EfficientNet-B3}
\resizebox{\columnwidth}{!}{
\begin{tabular}{lcccc}
\toprule
\textbf{Nhãn Mục tiêu} & \textbf{Precision} & \textbf{Recall} & \textbf{F1-score} & \textbf{Tổng (Support)} \\
\midrule
Bình Thường (Normal) & \textbf{0.91} & 0.94 & 0.92 & 27\\
Lành Tính (Benign) & \textbf{0.93} & 0.91 & 0.92 & 87\\
Ác Tính (Malignant) & 0.88 & \textbf{0.89} & 0.88 & 42\\
\midrule
\textbf{Tính tổng Macro Avg} & 0.90 & 0.91 & 0.90 & 156 \\
\textbf{Độ chính xác toàn cục} & - & - & \textbf{0.91(Accuracy)} & 156 \\
\bottomrule
\end{tabular}
}
\end{table}

\textbf{Phân tích độ chênh lệch:}
Tỉ lệ phát hiện Ác tính (Recall=0.89) là tương đối thành công vì cấu tạo bóng ác tính trên phim phản xạ chùm sóng (ultrasound echogenicity) hòa quyền 80\% hệt như cấu tạo tuyến sữa xung quanh thùy ngực bị vi vôi hóa. Một vài sai số (Âm tính giả Ác tính) xảy ra khi mô sẹo hậu phẫu thuật siêu âm hoàn toàn giống với khối u di căn xâm lấn rìa nham nhở gai tia. Trường hợp này AI sẽ bị nhầm FP thành Lành Tính (Fibroadenoma), tuy vậy về măt bác sĩ, hệ thống vẫn đang chỉ định "Cần Phẫu Thuật Kiểm Tra (Biopsy)", không thả bệnh nhân. Nên về mặt quy trình là được chấp thuận.

\subsection{Kết quả Thử nghiệm với ResNet-50 (Mô Sinh Thiết)}
Với bộ BreaKHis Test độc lập:
\begin{table}[H]
\centering
\caption{Chỉ số Đánh giá Mô hình ResNet-50 Dữ liệu Sinh thiết Tế bào Học (Histopathology)}
\resizebox{\columnwidth}{!}{
\begin{tabular}{lcccc}
\toprule
\textbf{Nhãn Phân Lớp} & \textbf{Precision} & \textbf{Recall} & \textbf{F1-score} \\
\midrule
Lành Tính (Benign Tế Khối) & 0.95 & 0.93 & 0.94 \\
Ác Tính (Malignant Ung Thư) & 0.96 & \textbf{0.97} & 0.96 \\
\midrule
\textbf{Chỉnh chung Macro Avg} & 0.95 & 0.95 & 0.95 \\
\bottomrule
\end{tabular}
}
\end{table}

Kết quả ở ResNet-50 cao cấp hơn siêu âm rất nhiều vì kính hiển vi chụp cấu trúc các biến dị nhân tế bài vô cùng rõ rệt: Tế bào ung thư cực kì dị dạng (pleomorphism), bào tương (cytoplasm) rộng đen màu, trong khi khối mào lành tính tạo thành cụm tế bào bao kín với màng đáy không xâm lấn. ResNet-50 với cơ chế Bottleneck layer phát hiện sự hỗn loạn của nhân tế bào ác tính không bị sai chệch bao nhiêu. F1-Score đạt 0.96 là kết quả hoàn hảo cho bài toán nhị phân.

\begin{figure}[H]
\centering
\vspace{0.5cm}
\includegraphics[width=\linewidth]{loss_curve.png}
\caption{Biểu đồ đường Learning Curves về Quá trình Hội tụ Mất Mát (Loss) của Mô hình Trọng Giao qua 25 Epochs (Khoảng trống dành riêng chèn Plot bằng Matplotlib).}
\end{figure}

\section{Trực quan hóa và AI mang tính Giải thích (Explainable AI - XAI)}
\subsection{Vấn đề Hộp đen Của Nơ - Ron (Block Box Limitation)}
Nút thắt cổ chai lớn nhất khiến Hội đồng y khoa (FDA Hoa Kỳ, hay Bộ Y Tế) ngần ngại khi sử dụng trí não nhân tạo ứng dụng khám bệnh không phải đoạt việc của họ, mà là "Liệu cái máy nó đang dựa vào Lỗ hổng (Artifact) nào để kết luận".
Ví dụ: Có thể trong tập Dữ liệu, Các bệnh nhân siêu âm u ác tính đều có dòng chữ "Tumor Scanner X2V" đóng watermark ở góc bản quyền. Do vô tình thiên kiến, Convolutional Layer đi học thuộc dòng chữ đó thay vì học khu vực khối u. Nếu AI chỉ cung cấp con số "Xác Suất 99\% Ác Tính", bác sĩ sẽ ngầm tưởng máy đang đọc lõi u, nhưng kỳ thực máy chỉ đếm pixels của dòng Watermark.
Để rọi ánh sáng toán học vào mạng lưới tối này, Phương pháp \textbf{Grad-CAM} (Gradient-weighted Class Activation Mapping) được tích hợp triển vọng vào luồng tương lai hệ thống.

\subsection{Công thức Thuật toán của Grad-CAM}
Grad-CAM \cite{gradcam} dựa vào Gradient (độ chênh lệch đạo hàm lùi) để chiếu lại khối ảnh gốc.
Gọi $y^c$ là giá trị logit dự đoán đầu ra (trước softmax) của nhóm nhãn ác tính mục tiêu $c$. Ta tìm các đặc trưng không gian đa chiều (Feature map) của chập Convolution Lõi Sâu cuối cùng, được đánh số là $A^k$ (Lớp filter thứ $k$).
Trọng số phân luồng kích hoạt toán học cho bộ $\alpha_k^c$ được tìm bằng cách sử dụng Global Average Pooling trên dòng chảy gradient:
\begin{equation}
\alpha_k^c = \overbrace{\frac{1}{Z} \sum_{i} \sum_{j}}^{\text{Global Average Pooling}} \underbrace{\frac{\partial y^c}{\partial A_{ij}^k}}_{\text{Gradients by Backprop}}
\end{equation}
Bản đồ chói lọi (Heatmap localization map) $L_{Grad-CAM}^c$ 2 Chiều sẽ là tổng dồn tuyến tính của hàm kết hợp (Linear Combination) sau đó khóa qua cơ chế lọc biến phi tuyến dương ReLU:
\begin{equation}
L_{Grad-CAM}^c = ReLU \Big( \sum_k \alpha_k^c A^k \Big)
\end{equation}
Lý do bắt buộc phải sử dụng hàm kích hoạt ReLU ở đây là vì ta chỉ quan tâm đến những vùng không gian hình ảnh có tác động \textbf{tích cực dương} (Positive Influence) trong việc xúi giục AI kết luận bệnh Ác Tính. Những vùng có trọng số đạo hàm âm thường là đặc trưng nhận diện của lớp từ chối (Ví dụ: Đặc trưng của Lành Tính), nên việc giữ lại chúng trên bản đồ nhiệt sẽ gây nhiễu loạn trầm trọng thị giác của y bác sĩ.

Hàm Gradient sau phép chiếu này trả về 1 tensor mảng nhỏ $7 \times 7$ (Kích thước gốc của lớp chập cuối cùng từ ResNet-50). Vi xử lý nội bộ tại Server backend sẽ dùng kỹ thuật dãn nở ma trận hình học \textbf{Bilinear Interpolation Up-sampling} để phóng to kích thước Vector này về lại khuôn tỷ lệ gốc của hình ảnh bệnh nhân ban đầu ($500 \times 500$). Lớp ma trận nhiệt đới (Jet-Colormap RGB) sau đó được nội suy pha trộn đè trực tiếp lên bề mặt ảnh siêu âm xám gốc với phương trình cộng trọng số ngầm định Alpha ($\alpha = 0.5$):
\begin{equation}
\text{Image}_{Overlay} = \alpha \cdot \text{Image}_{Original} + (1-\alpha) \cdot \text{Map}_{Heat}
\end{equation}
Qua quá trình đồng bộ hóa đa tầng đó, vùng nào bức xạ Đỏ rực (Red Hot), chứng tỏ thuật toán Deep Learning đang dồn toàn lực rọi "ánh hào quang" nhận thức vào vùng tế bào bất định cấu trúc ấy để chốt sổ chẩn đoán lâm sàng cuối cùng.

\section{Hệ thống Ứng dụng Backend / REST Engine Design}
Luận lý chuyên phân tích toán học không thể thành công nếu nó chỉ nằm trên Terminal mã Code Python. Một phần lớn báo cáo phân bổ thiết kế khối kỹ sư kiến trúc phần mềm tích hợp (Software Architecture and Microservice).
Flask Framework siêu gọn của ngôn ngữ Python đã được tận dụng xây dựng Gateway. Mô hình WSGI (Web Server Gateway Interface) cho phép nhận đa luồng hình ảnh.

\subsection{Mô hình luồng dữ liệu theo thời gian thực (Micro-Inference System)}
Lõi ứng dụng (`app.py`) có khả năng giữ bộ nhớ mô hình PyTorch (Khoảng 200 Megabytes hằng trăm số Float32) tải chết trên lượng RAM (In-Memory Processing).
\begin{lstlisting}[language=Python, caption=Cấu trúc API Logic phân luồng ảnh Y khoa Server-side]
@app.route('/predict/ultrasound', methods=['POST'])
def handle_ultrasound():
if 'file' not in request.files:
return jsonify({"error": "No file part"}), 400
file = request.files['file']
# Doc bit-stream cua file qua bo dem buffer RAM roi day di inference thay vi mat I/O load ra SSD roi load vao
# Tra nguyen the ket qua % confidence dang mang JSON
result = predict_ultrasound(file.read())
return jsonify(result), 200
\end{lstlisting}

Kiện toàn hóa API: Bất kì ứng dụng di động app React Native, hay Web ReactJS, VueJS hay Web HTML thuần (như code ở `./frontend/index.html`) chỉ việc bắn POST Request Multipart/Form-Data là kết quả nhận về luôn dưới 0.3 Giây tốc độ trễ.
Mẫu hồi đáp Payload chuẩn:
\begin{lstlisting}
{
success: true,
prediction: "Malignant",
confidence: 98.45,
details: {
Malignant: 98.45,
Benign: 1.25,
Normal: 0.30
}
}
\end{lstlisting}

\subsection{Sơ đồ luồng giao diện Web Frontend trực quan}
File trình duyệt frontend giao kết bằng Javascript thuần túy thông qua API fetch. Bác sĩ (End-user) trải nghiệm mượt mà không cần tái phân phối trang Web (Single Page App Behavior). Kết xuất ra vạch trượt Gauge-bar đo mức lưu lượng phần trăm. Báo động nhãn màu Đỏ sậm nếu khối u được máy gán là nguy cấp sinh mệnh chuẩn bị báo cáo ca mổ khẩn cấp.

\subsection{Kết quả Giao diện Phần mềm và Giải thích AI (XAI)}
Thông qua quá trình tích hợp sâu PyTorch Grad-CAM vào pipeline API inference, hệ thống hiển thị thành công bản đồ nhiệt tự động từ lõi Convolutional lên giao diện phản hồi theo thời gian thực (Real-time Feedback).
\begin{figure}[H]
\centering
\includegraphics[width=\linewidth]{Screenshot 2026-03-04 011054.png}
\caption{Giao diện Hệ thống chẩn đoán thực tế, bản đồ nhiệt Grad-CAM định vị chính xác vị trí tâm lõi khối u ác tính so với viền cơ xương.}
\end{figure}

\section{Đánh giá, Hạn chế Rủi ro và Tương lai}
\subsection{Tổng kết}
Các thiết lập thực nghiệm sâu sắc này kết luận được Hệ thống Trí Báo Nhân Tạo nhận diện tế bào và các nốt đốm hình ảnh khối u vượt ngoài dự kiến đối với khả năng vi chỉnh tự tri thức bằng giải thuật tối ưu của Transfer Learning (Đặc biệt EfficientNet Compound và Resnet Vòng cung lõm). Với mô đun độc quyền chia đôi phương thức phục vụ linh hoạt cho đa phương tiện y tế, sản phẩm đem lại quy trình tích hợp khép kín không phụ thuộc bên 3rd Cloud Platform nào.

\subsection{Khuyến cáo đạo đức về AI Thuật Toán Trợ Y}
Tất cả những con số AI dự đoán ra MẶC ĐỊNH LÀ KHÔNG ĐƯỢC PHÁP LUẬT Y TẾ THỪA NHẬN dùng làm cơ sở viết bệnh án phác đồ kê đơn truyền hóa chất (Chemo) cắt vú. Bản thân mô hình dù Acc 98\%, số lượng $2\%$ sai số mạng nhầm FP FN là đại diện cho hàng vạn sinh mạng vô tội. Do vậy mô hình sinh ra trong đồ án, hoặc ngay cả bệnh viện đa khoa trung ương luôn mang mác "Hệ Thống Trợ Lý - CAD". Quyết định sinh sát cuổi cùng phải là nhóm Chuyên Gia Y Khoa Tổng Hợp Hội chẩn (Doctor-In-The-Loop Hệ Thống Có Nhân Sự Theo Dõi).

\subsection{Tương lai triển vọng Nghiên cứu Vĩ mô và Quản trị Hệ Sinh thái Y tế Số}
Mặc dù đang ở đỉnh cao thiết kế mô phỏng (In-Silico Design), các mô hình học sâu y sinh luôn rình rập những rào cản ngăn giới khoa học tự đắc. Để hệ thống nội bộ của Đồ án có đủ tư cách thương mại hóa (Commercialization) và lắp đặt thực tiễn tại bệnh viện tuyến tỉnh, các luồng kiến trúc sau bắt buộc phải được thiết kế và vá lỗi trong tương lai:

\begin{itemize}
\item \textbf{Dịch chuyển Tiền miền dữ liệu Không gian (Domain Shift \& Generalization)}: Trở ngại tử huyệt của Mạng lưới AI Y tế hiện đại là hội chứng "Học Thuộc Lòng Thiết bị" (Scanner Hardware Overfitting). Một đồ án được huấn luyện cực tốt, F1-Score lên tới 99\% trên tệp máy quét Siêu âm đời cũ của bệnh viện Baheya (Ai Cập) có thể đột ngột giảm sút độ chính xác thê thảm xuống 60\% khi đi cài đặt thực tế trên máy quét Philips Epiq đời mới của Bệnh viện K (Việt Nam). Sự sai khác về dải phổ tần số tần âm, thuật toán làm mịn bộ lọc phần cứng (Hardware artifact) vô tình bẫy AI nhận lầm mầm bệnh. Việc mở rộng cơ chế Thích ứng Miền Thích Nghi sâu (Deep Domain Adaptation) ép buộc mô hình phải học các \textit{Đặc trưng Bất biến tính Môi trường (Environment-Invariant Features)} thay vì bị ám ảnh bởi độ mịn phân giải của một cỗ máy duy nhất là nhiệm vụ sống còn.
\item \textbf{Bảo mật Học Máy Liên kết Máy chủ Trung tâm (Federated Learning Security)}: Định luật Bảo vệ dữ liệu cá nhân Y tế trên quy mô toàn cầu cực kỳ khắt khe. Đặc biệt là bộ luật HIPAA của chính phủ liên bang Hoa Kỳ hay khối GDPR của Châu Âu, nghiêm cấm hình sự hóa việc các bệnh viện lưu trữ copy-paste hàng triệu thẻ bệnh án nhạy cảm của bệnh nhân qua email hoặc đẩy tự do lên chung một vùng Cloud nội trú chắp vá để Train thuật toán. Công nghệ Phi tập trung Liên Kết (Federated) ra đời cho phép mỗi Bệnh viện Trung ương trên thế giới tự tải bản Trọng số Thô của mô hình đồ án này về Train hoàn toàn cục bộ trên kho máy chủ sinh thiết mật của bệnh viện đó. Sau khi hội tụ, thuật toán rễ máy khách chỉ \textbf{Tính toán và gửi Ngược Băng thông dãy sô Đạo Hàm Gradient Deltas} lên Data-Center Máy chủ Tổng của đồ án này để hòa quyện tri thức với Trọng số chung (Global Aggregation), mà không hề rò rỉ bất cứ một hạt pixel tĩnh nào liên quan danh tính khuôn mặt, mã định danh hay thân thể ngực nhũ quan của người phụ nữ ra Internet ngoài.
\item \textbf{Cơ chế Phân tách Khung lưới Tọa độ Kép (Bounding-Box Semantic Segmentation)}: Để thay thế việc màn hình máy chỉ nhả duy nhất kết luận cụt lủn "Tỷ lệ chẩn đoán bức ảnh này mang nhãn Malignant là 98\%", Đồ án cần được tài trợ mở rộng nghiên cứu bằng việc tích hợp Mạng cấu trúc Hình ống khói chữ U cực thịnh hành - \textbf{U-Net} (Đoạt giải Nobel cho khối ngành Tích chập). Khởi τrạy U-Net trên tệp dữ liệu viền Mặt nạ tinh khiết (Ground-Truth Masks Dataset) để Auto-Lắng Hình, máy tính sẽ trực tiếp tô Vẽ Khối đa giác vùng bóng trắng bao bọc chuẩn xác tới milimet vuông xung quanh lõi viền bướu cho mắt thường của bác sỹ hội chẩn theo dõi. Thuật toán sau ranh giới Segment hóa còn có thể thiết lập phương trình đo trực tiếp Kích thước tiết diện bề mặt Cắt ngang 2D ($cm^2$), tham số vàng để in ra hồ sơ theo dõi khách quan tốc độ di căn phình to / thu cạn của khổi u theo từng chu kỳ thuốc hóa trị.
\item \textbf{Cấy Rễ Cơ thể Cơ chế Nhận thức Thị giác Trọng tâm (Vision Transformers - ViT)}: Tương tự như sức mạnh nhảy vọt đột phá của ChatGPT và Gemini trong xử lý Ngôn ngữ Tự Nhiên (LLM NLP), việc ứng dụng Cấu trúc \textit{Trường tập trung Tự Chú Ý Định tuyến (Self-Attention Mechanism)} nguyên bản đè lên Dữ liệu Hình Ảnh sẽ giúp AI tự tạo bầu trời ma trận kết nối đa cực. ViT tự học luật liên hệ tầm xa (Long-range dependencies) giữa tận mút viền ngoài cùng của cấu trúc khối u với nhân đậm trung tâm của nó trong một nhịp xử lý đồng bộ, đè bẹp hoàn toàn kỷ nguyên trượt khung chập cửa sổ thiển cận lặp lại $3 \times 3$ của nền hành chính CNN truyền thống cổ điển.
\item \textbf{Kiến trúc Cảng Mạng Đám Mây Chuyên biệt SaaS (Software-as-a-Service Framework)}: Gói gọn đồ án tại mức tinh cầu cục bộ (Localhost) thành một tệp nhị phân Python độc lập Container hóa (Docker Image Sandbox). Quản lý chuỗi luân chuyển phân luồng Worker bằng Gunicorn HTTP Server, kết hợp hệ thiết bị Phân chia Tải (Traefik Load Balancer) để vận hành lơ lửng trên hạ tầng kết hợp mây dịch vụ AWS (Amazon Web Services). Điện toán ngầm cấu hình Máy chủ gắn Chip gia tốc phần cứng Nvidia Tensor Core \textbf{CUDA Inference GPUs}. Thiết lập đường dẫn này giúp hội đồng nhóm nghiên cứu đúc kết thành Công ty thương mại phần mềm đóng gói, bán cấp phép định tuyến API đường truyền dữ liệu sinh học rẻ mạt với trễ cực thấp đến mọi phân viện lâm sàng quy mô vi mô cực nhỏ chìm dưới các vùng hải đảo núi cao, tận cùng biên giới hẻo lánh của tổ quốc - những nơi ngập rủi ro, thiếu thốn hoàn toàn thiết bị tầm soát Mammography y khoa tối tân, mà họ chỉ mang độc một sóng truyền tín hiệu mạng mỏng manh 4G Internet qua điện thoại vệ tinh.
\end{itemize}

\section{Kết luận Toàn Văn}
Trong kỷ nguyên Chuyển đổi số Y tế Công cộng hiện nay, thuật toán Trí thông minh là bánh răng không thể trễ trễ ở hệ thống lâm sàng đa dạng. Báo cáo Tốt nghiệp Đồ án AI này đã đi vòng quanh 10 chương mục liên tục mô phỏng phương pháp biến lý thuyết của các nhà khoa học Google/Facebook (Pytorch Framework, CNN EfficientNet) xuống thẳng thực hành một trang Web Y Tế Phân Lớp nhanh chuẩn, giá rẻ vận hành bằng Máy Đơn, nhưng mang năng lực Tầm Cỡ Toàn Cầu. Tương lai ngập đầy hứa hẹn, hướng đích giảm tải bác sĩ, đưa máy dự đoán phòng ung thư đại chúng cho đồng bào mọi lúc trạm xá, gỡ phán quyết sai phạm nhũ quan.

\begin{thebibliography}{00}
\bibitem{busi} Walid Al-Dhabyani, M. Gomaa, H. Khaled, A. Fahmy, \emph{Dataset of breast ultrasound images}, Data in Brief, vol. 28, 2020, 104863.
\bibitem{breakhis} F. A. Spanhol, L. S. Oliveira, C. Petitjean, L. Heutte, \emph{A Dataset for Breast Cancer Histopathological Image Classification}, IEEE Transactions on Biomedical Engineering, vol. 63, no. 7, pp. 1455-1462, 2016.
\bibitem{resnet} K. He, X. Zhang, S. Ren, and J. Sun, \emph{Deep Residual Learning for Image Recognition}, in IEEE Conference on Computer Vision and Pattern Recognition (CVPR), 2016, pp. 770-778.
\bibitem{efficientnet} M. Tan and Q. V. Le, \emph{EfficientNet: Rethinking Model Scaling for Convolutional Neural Networks}, in International Conference on Machine Learning (ICML), 2019, pp. 6105-6114.
\bibitem{gradcam} R. R. Selvaraju, M. Cogswell, A. Das, R. Vedantam, D. Parikh, and D. Batra, \emph{Grad-CAM: Visual Explanations from Deep Networks via Gradient-Based Localization}, International Journal of Computer Vision, vol. 128, pp. 336-359, 2020.
\bibitem{alexnet} A. Krizhevsky, I. Sutskever, and G. E. Hinton, \emph{ImageNet Classification with Deep Convolutional Neural Networks}, Advances in Neural Information Processing Systems, vol. 25, 2012.
\bibitem{adam} D. P. Kingma and J. Ba, \emph{Adam: A Method for Stochastic Optimization}, International Conference on Learning Representations (ICLR), 2015.
\bibitem{pytorch} A. Paszke, et al., \emph{PyTorch: An Imperative Style, High-Performance Deep Learning Library}, Advances in Neural Information Processing Systems, pp. 8024-8035, 2019.
\bibitem{cad_review} L. Wang, T. K. Tan, R. C. Fernandes, \emph{Computer-Aided Diagnosis System for Breast Cancer}, IEEE Reviews in Biomedical Engineering, vol. 10, pp. 15-30, 2017.
\bibitem{data_aug} L. Perez, J. Wang, \emph{The Effectiveness of Data Augmentation in Image Classification using Deep Learning}, arXiv preprint, 2017.

\end{thebibliography}

\end{document}